\chapter{绪论}
\label{chap:1}
\chapterguide{说明研究动机、已有方法、研究范围和论文主线。研究现状需用正式阅读的文献补齐，不能把项目设计直接写成已证明的学术创新。}

\section{研究背景与意义}
\label{sec:1-1}

\subsection{异构CAD环境下的模型交换需求}
\label{subsec:1-1-1}
\outlinehint{从协同设计、软件迁移和历史零件复用引出跨系统交换需求，结合实际工作场景说明目标模型仍需继续修改。}

\subsection{参数化模型转换中的信息保留问题}
\label{subsec:1-1-2}
\outlinehint{区分外形、特征类型、尺寸参数、引用关系和修改后的更新行为；用定深切除与到面切除的区别说明仅外形一致仍不充分。}

\subsection{可编辑特征转换的研究意义}
\label{subsec:1-1-3}
\outlinehint{围绕减少重复建模和支持模型延续使用阐述意义。零件库检索、批量变体与训练数据只作为潜在用途，不作为已验证成果。}

\section{国内外研究现状}
\label{sec:1-2}

\subsection{几何与拓扑信息交换方法}
\label{subsec:1-2-1}
\outlinehint{调研中性格式和几何内核交换的实际信息范围，区分本项目使用的终态几何导出与标准本身可能支持的更广信息。}

\subsection{基于特征回放与宏参数化的转换方法}
\label{subsec:1-2-2}
\outlinehint{比较特征序列、参数和约束的提取与重建机制，关注映射范围、源目标依赖和参考对象处理方式。}

\subsection{基于统一模型的特征转换方法}
\label{subsec:1-2-3}
\outlinehint{比较中间表达与语义模型如何组织参数和依赖，讨论系统扩展、参数差异与目标重建的研究问题。}

\subsection{转换结果验证方法与现有研究不足}
\label{subsec:1-2-4}
\outlinehint{分别梳理结构、参数、几何和编辑验证方法，在真实文献对照后归纳不足，不把几何成功当作可编辑性已完整保留。}

\section{研究目标与主要研究内容}
\label{sec:1-3}

\subsection{形变特征定义与研究范围}
\label{subsec:1-3-1}
\outlinehint{沿用中期题目中的“形变特征”，明确指基于草图与参数驱动的成形特征，不指自由曲面或网格变形。以拉伸、旋转、扫掠为核心，以草图、基准面为支撑；轴、点、边、面仅作为相关引用涉及。}

\subsection{研究目标与关键问题}
\label{subsec:1-3-2}
\outlinehint{提出四个问题：参数实际含义如何统一，空间与引用如何恢复，公共表达如何转为目标建模步骤，以及保留程度如何评价。}

\subsection{主要研究任务与个人工作边界}
\label{subsec:1-3-3}
\outlinehint{组织差异分析、统一模型、转换算法、系统实现与实验任务；按实际分工说明扫掠的主导与协作部分，其他工程特征不列为个人核心贡献。}

\section{技术路线与论文结构}
\label{sec:1-4}

\subsection{总体技术路线}
\label{subsec:1-4-1}
\outlinehint{给出“几何构造分析—三系统参数调研—公共表达—目标重建—分层验证”的路线，并标明实验反馈如何修正字段和映射。}

\subsection{论文章节安排}
\label{subsec:1-4-2}
\outlinehint{说明第二章回答差异和需求，第三章定义保存什么，第四章解释怎样转换，第五章说明系统落地，第六章评价结果，第七章收敛结论。}
